\pdfminorversion=7
\documentclass{article}
\usepackage[dblblindworkshop,final]{neurips_2026}
\workshoptitle{ML4Molecules: Agentic Systems for Molecular Sciences}
\usepackage[utf8]{inputenc}
\usepackage[T1]{fontenc}
% Latin Modern typewriter: the same design as the default EC typewriter face,
% but as Type 1 fonts rather than bitmaps (cm-super is not part of a base TeX Live).
\renewcommand{\ttdefault}{lmtt}
\usepackage{amsmath,amssymb,bm}
\usepackage{booktabs}
\usepackage{tabularx}
\usepackage{float}
\usepackage{graphicx}
\usepackage{microtype}
\usepackage[hidelinks]{hyperref}
\usepackage{url}
\usepackage{xcolor}
\usepackage{multirow}
\newcommand{\vect}[1]{\bm{#1}}
\newcommand{\method}{ROSHAMBO3}
\input{generated/constants.tex}
\input{generated/development-constants.tex}
\input{generated/roofline-constants.tex}
\input{generated/constants-optimizer-sensitivity.tex}
\input{generated/portability-constants.tex}
\input{generated/cpu-scaling-constants.tex}
\graphicspath{{generated/figures/}}

\title{Linear-time Shape Alignment with ROSHAMBO3}

\author{Imran S. Haque\\
Kimia Therapeutics, Berkeley, CA 94710\\
\texttt{ihaque@kimiatx.com}}

\begin{document}
\maketitle
\vspace{-18pt}
\begin{abstract}
\vspace{-2pt}
    Three-dimensional (shape) similarity is a useful primitive for virtual screening
    and molecular machine learning, but current open implementations are limited in performance
    by a quadratic inner loop across pairs of query- and database-molecule atoms.
    \method{} achieves a \VoltaShapeCorrectedSpeedupWhole--\VoltaComboStockSpeedupWhole$\times$ speedup
    over the previous open SOTA ROSHAMBO2 by using query-amortized Gaussian fields that make the
    inner loop linear combined with system-scale optimization
    at the algorithmic, GPU, and host layers, screening up to \DfiveMillionEvaluationsPerSecondWhole{}M query--conformer
    pairs/s on NVIDIA GPUs, with secondary CPU (x86, ARM) and Apple Metal backends.
    We further describe the human expert--AI interaction loop that guided the exploration and
    implementation of \method{} and discuss implications for the scientific software commons.
\end{abstract}

\section{Introduction}

Ligand-based virtual screening searches a library for compounds similar to a
known query. Methods broadly divide into ``two-dimensional'' comparisons of
molecular graphs and ``three-dimensional'' comparisons of conformations.
Although 2D fingerprints are
inexpensive and effective, they intrinsically bias chemical search towards
regions of \emph{chemical} similarity, rather than \emph{property} similarity.
Molecular properties like biological interactions are often rooted in
steric and electrostatic complementarity; consequently 3D similarity
enables ``scaffold hopping'': the retrieval of topologically-distinct molecules
bearing compatible volume and interaction patterns~\citep{rush2005scaffold}.

Three-dimensional search follows two main approaches. Embedding methods compare
molecules in an abstract vector space
\citep{haque2010scissors,schreyer2012usrcat,atsango2022molclass,winter2026sand};
direct-overlap methods instead optimize a volumetric kernel over rigid-body
transformations in $SE(3)$. Embeddings can be very fast but are less interpretable, as they do not compute visualizable poses.
Search over an enumerated database can also discretize an abstract desired 3D
shape into a physical compound, making fast direct search useful both in drug
discovery and as an action or reward function for chemical-design agents.

Gaussian-overlap methods represent atoms as smooth volumes~\citep{grant1995gaussian,grant1996fast}; their proprietary implementation in the OpenEye ROCS package is widespread.
PAPER~\citep{haque2010paper} first mapped this calculation to GPUs, followed by OpenEye's FastROCS package. Subsequently ROSHAMBO and ROSHAMBO2
provided open implementations modernizing and extending shape similarity to pharmacophore-based ``color'' similarity~\citep{atwi2024roshambo,atwi2025roshambo2}.

We present ROSHAMBO3, which accelerates ROSHAMBO2 by \VoltaShapeCorrectedSpeedupWhole--\VoltaComboStockSpeedupWhole$\times$,
enabling one query to search \ThreeCrMillionConformersWhole{} million conformers in \ThreeCrQOneWallSecondsWhole~s on one GPU.
The numeral is not just a version, but also denotes ``THREE'' (\emph{Tokens Help Researchers Expedite
Engineering}): a human expert built an operating environment for
Anthropic's Claude Code and OpenAI's Codex agents to explore the design landscape with periodic feedback, illustrating potential gains in the
scientific software commons through modest investments in expert-AI collaboration.
ROSHAMBO3 is open source at \url{https://github.com/kimiatx-oss/roshambo3}; the
analyses, run records, and an export of the development record behind this paper are at
\url{https://github.com/kimiatx-oss/roshambo3-paper}.




\section{Methods}
\label{sec:method}
\paragraph{Shape and color objectives.}
For query atoms $i\in A$ and candidate atoms $j\in B$ transformed through rotation $R$ and translation $t$, the
``analytic'' first-order Gaussian cross-overlap has the form of a quadratic-complexity summation over
atom pairs, with the commonly used Tanimoto coefficient $T$ defined as a simple function of 
overlap and molecule self-volumes:
\begin{equation}
 V_{AB}(R,\vect t)=\sum_{ij} C_{ij}
 \exp[-\beta_{ij}\|\vect a_i-(R\vect b_j+\vect t)\|^2],
 \quad T=\frac{V_{AB}}{V_{AA}+V_{BB}-V_{AB}} .
 \label{eq:overlap}
\end{equation}

This function extends to multiple atom \emph{planes} with defined inter-plane
weights. ``Color'' adds virtual atoms at pharmacophore centroids, such as hydrophobes,
H-bond donors or acceptors, aromatic rings, and ionizable groups. In practice,
the plane interaction matrix is diagonal: shape and each color type interact only
with themselves, with one additional parameter weighting shape against summed color. Adding
color also admits distinct objectives: shape-only, ``combo'' (shape and color
optimized at each step), and ``rescore'' (pose optimized by shape then rescored with shape and color
at the final pose).



\begin{figure}[t]
 \centering
 \includegraphics[width=\textwidth]{field-schematic.pdf}
 \caption{\textbf{A:} Analytic overlap evaluates all query--library atom pairs
 per step. \textbf{B:} \method{} rasterizes the query once; each transformed
 library atom samples the field value and gradient.}
 \label{fig:field}
\end{figure}

\paragraph{Linear time overlap with query fields.}
Search is asymmetric: one query is reused for many candidates. Since Gaussian overlap
is linearly additive, it is possible to precompute a discretized \emph{field}-based
representation of the query $A$, making the overlap computation a single linear-time
loop over library molecule atoms (Figure~\ref{fig:field}):
\begin{equation}
 \Phi_A(\vect x)=\sum_i C_i e^{-\beta_i\|\vect x-\vect a_i\|^2},
 \qquad V_{AB}(R,\vect t)=\sum_j w_j\Phi_A(R\vect b_j+\vect t).
 \label{eq:field}
\end{equation}

The $x/y/z$ gradient fields of the query can be similarly discretized to avoid
first-difference gradient discontinuities that can trap optimizers. This representation
is a natural fit for GPUs, where texture units perform hardware acceleration of spatially-local
interpolated lookups, enabling fast trilinearly-interpolated lookup of
both value and gradients simultaneously.

Similar transformations with grid alignment and fast Gaussian sums
have substantial prior art \citep{ronkko2006brutus,tsin2004correlation,gao2019filterreg,greengard1991fast},
and OpenEye documentation exposes, but does not elaborate on, a ``Grid'' method \citep{openeyeShapeTheory}.
Surprisingly, we have found no prior openly documented implementation
exploiting this transformation and interpolating both value and gradient for molecular shape.

\paragraph{Decoupled rigid-body Adagrad.}
Prior implementations of Gaussian shape optimization represented the $SE(3)$ rigid
body transformation as a unit quaternion rotation with a translation and optimized this 7-vector using BFGS or Adagrad \citep{duchi2011adaptive}.
Both optimizers face two representational liabilities. 
Unit quaternions representing
rotations form a 3-dimensional manifold \(S^3\) embedded in \(\mathbb{R}^4\), but the optimizers
do not know about this structure. Their Euclidean steps, which both require renormalization and have non-constant angular velocity depending on starting point, make rotational reach depend on
the current quaternion and create poorly conditioned regions. Secondly, the rotational and
translational components of the rigid-body transformation share neither units nor scale and so may not benefit from sharing an optimizer velocity.



ROSHAMBO3's optimizer solves both issues. For the first,
we rotate by tangent-space multiplication:
\begin{equation}
 q_{k+1}=\exp(\tfrac12\delta\vect\theta_k)\otimes q_k,
 \quad \delta\vect\theta_k=\eta_q\vect g_{q,k}/
 \sqrt{\vect G_{q,k}+\epsilon},
 \label{eq:expmap}
\end{equation}
For the second, we apply Adagrad separately
for rotations and for translations with separate accumulators and learning rates.
A particular benefit is that a larger translation learning rate
enabled larger reach in optimization (the translation travel budget scales as $\eta_t\sum_k k^{-1/2}$),
enabling ROSHAMBO3 to reduce optimizer steps to 25 from ROSHAMBO2's 100 without shrinking the search basin.
Because the search itself changed, rankings do not transfer exactly between engines: over \OptSensQueries{} PubChem queries, molecule-level top-\OptSensTopK{} lists from ROSHAMBO3 and correctness-fixed ROSHAMBO2 overlap by \OptSensComboTopKOverlapPct\% (combo) and \OptSensShapeTopKOverlapPct\% (shape).
Rescoring each engine's final poses with an independent float64 reference reproduces its own scores, so the difference lies in which optima are found rather than in scoring; on combo, \OptSensComboDisplacedPct\% of molecules leaving ROSHAMBO2's top-\OptSensTopK{} were displaced by molecules ROSHAMBO3 optimized to higher scores, not scored lower themselves.


\paragraph{GPU execution and resource balance.}
For a primer on GPU architecture, see Appendix Section~\ref{sec:gpumodel}.
Shape similarity without early stopping has a deterministic operation count for any given count of query and database atoms, and the operation mix differs for analytic and field kernels. We used this fact to compute roofline estimates for each kernel, modeling the maximum possible achievable throughput for a kernel assuming it were limited by a single resource, based on hardware specifications,
corpus statistics, profiled instruction mixes, and observed power-limited clocks
(Appendix~\ref{app:roofline}, Table~\ref{tab:roofline-device-caps}). Fields move the bottleneck away from special function (exponential) evaluation, which caps the analytic kernel, and substantially raise the ceiling, now set by different functional units on different GPUs.
GB202 rises from \RooflineBlackwellAnalyticStructuralMillionPairsPerSecond{}M to \RooflineBlackwellFieldIssueMillionPairsPerSecond{}M pairs/s, limited by instruction issue.
AD104 (L4) and GV100 rise from \RooflineLFourAnalyticStructuralMillionPairsPerSecond{}M and \RooflineTitanAnalyticStructuralMillionPairsPerSecond{}M to \RooflineLFourFieldLoneTexMillionPairsPerSecond{}M and \RooflineTitanFieldLoneTexMillionPairsPerSecond{}M pairs/s respectively, both limited by L1TEX (texture interpolation) rate; the GV100 limit is projected from an assumed \RooflineTitanLoneTexSectorsPerSmCycle{} sectors/SM/cycle, close to the \RooflineLFourLoneTexSectorsPerSmCycle{} measured on AD104.

Simply increasing the roofline does not guarantee improved performance; to maximize the efficiency of ROSHAMBO3 kernels we performed extensive testing of configurations including atom and starting-pose parallelism, coordinate placement (registers, shared memory, or cached global memory), L1/shared-memory carveout, cache hints, and whether color planes used field or analytic evaluation. No single configuration dominated in every scenario, especially between shape and combo optimization. The shipped configurations of ROSHAMBO3 are documented in Appendix~\ref{app:kernel-config}.


\paragraph{Precision, storage, and the streaming pipeline.}
Maximizing GPU throughput also requires optimization of the producer (upload of data to the GPU) and consumer (reduction and result readback) pipelines. ROSHAMBO3 introduces a new on-disk format optimized for throughput: coordinates are quantized to a $2^{\CoordinateStepLogTwo}$~\AA{} int16 grid and bit-packed per block at the narrowest width each axis needs, bucketed by atom count to minimize warp padding waste, with color planes ordered by measured pair work to maximize bounded cache effectiveness.
A pool of CPU producer threads decodes independent blocks natively into the kernel layout and fills pinned-memory buffers. Field textures
 are stored in fp16 format as a \texttt{half4} texture hardware-expanded to fp32 for computation.
Finally, a single GPU-facing
consumer preserves corpus order while overlapping copies with earlier kernels,
caching query fields, and coalescing small chunks into larger launches. This
pipeline became necessary because acceleration of field evaluation exposed
host reduction, rather than the kernel, as the end-to-end bottleneck. The
current consumer performs a deterministic molecule-level top-$k$ selection on
the device and returns only the bounded hit list. An earlier design returned
every conformer score over PCIe; eliminating this device-to-host stream was
as important to system performance as accelerating the scoring
kernel.

\paragraph{Interfaces and backends.} One engine is exposed through a Python API,
a command-line tool (\texttt{roshambo3 prepare} builds a corpus,
\texttt{roshambo3 screen} searches it), and a lab web application with SMILES
or file queries, 3D pose overlays, and CSV/SDF export. Each runs on the CUDA
backend, a native x86/ARM CPU backend, or an experimental Metal backend.

\clearpage
\paragraph{Agentic exploration.} ROSHAMBO3 started as chat conversations
between the author and Google's Gemini 3.0 Pro Preview about how to
accelerate shape similarity with precomputed fields or new hardware units on modern GPUs. Implementation
and further exploration were almost entirely carried out by AI agents in two agent systems sharing one repository --- Anthropic's Claude Code (Claude Sonnet and Opus 4.8 through 5.5, and Fable 5) and OpenAI's Codex (GPT~5.6-Terra/Sol, then Sol, Astra, and Luna 6 and Sol and Luna 6.1) --- with higher-level feedback from the author; Gemini 3.1 Pro Preview also reviewed early work. To manage parallel agent swarms, a discipline was maintained with well-documented tickets and plans checked into a defined location in the repository, and agent protocols around ticket checkout and evaluation criteria. Specifically, each run was required to record source, build, GPU, clock and power telemetry, cache state, and repeat phase, with exclusive locks over GPUs to prevent
parallel agents from interfering with each others' runs. Bitwise identity
tests were required to pass for changes that should not change arithmetic
(eg, memory staging), with a retrieval gate (top-K identity) used for changes inducing subtle (FP associativity) or more drastic (field vs analytic, fp16/32 quantization, optimizer) numeric changes.

\section{Evaluation}
\label{sec:evaluation}

\paragraph{Datasets.} 
To assess correctness of ROSHAMBO3, we used the LIT-PCBA and DUDE-Z datasets
on which ROSHAMBO2 was evaluated~\citep{trannguyen2020litpcba,mysinger2012dude}.
We requested 50 conformers per molecule using ETKDGv3~\citep{riniker2015etkdg}
followed by MMFF94s minimization, producing
\LitpcbaCorpusConformers{} conformers for LIT-PCBA (\LitpcbaDiskGigabytes~GB).
PANEL-8 is a deterministic set of one deposited ligand from each of eight target/assay
groups; its exact PDB coordinates are retained and span
\LitpcbaQueryHeavyAtomRange{} heavy atoms. Thus, the public timing arm contains
realistic shape and color records; enrichment results are included to
establish parity rather than to demonstrate the utility of shape screening
itself.

To assess ultrascale performance, we benchmarked on two large datasets. ``PubChem3D''
is a shape-only build of PubChem3D~\citep{bolton2011pubchem3d} containing \DfiveMolecules{} molecules and
\DfiveConformers{} conformers (mean \DfiveMeanConformersPerMolecule{} conformers/molecule, \DfiveDiskGigabytes~GB on disk). As a real-world test case we also benchmark
search inside a Kimia-proprietary virtual hit-finding library containing 75.2M molecules (mean \ThreeCrMeanConformersPerMolecule{} conformers/molecule; \ThreeCrDiskGigabytes~GB) with color. PubChem3D
was queried with the first four PANEL-8 structures; the proprietary library was
queried with four conformers of a single \ThreeCrQueryHeavyAtoms-heavy-atom private molecule
relevant to an internal discovery program.

\paragraph{Hardware.} We benchmark on seven machines (Appendix Table~\ref{tab:systems}): workstations
with an NVIDIA RTX PRO 6000 Blackwell Max-Q (``GB202'') and a TITAN V (``GV100''), an M4 Pro MacBook Pro (14-inch, Nov 2024),
and four VMs on Google Cloud: one GPU host with an NVIDIA L4 (g2-standard-16), and three
CPU-only VMs (c4/c4d/c4a-highcpu-16: Intel Xeon 8581C ``Emerald Rapids'', AMD EPYC 9B45 ``Turin'', and ARM Neoverse V2 ``Google Axion'', respectively).

\paragraph{Accuracy and performance relative to ROSHAMBO2.} Gaussian shape
overlay is intrinsically approximate: infinite-extent atoms are unphysical, and
first-order overlap overcounts volume. Rather than anchoring on exact computation of an approximate value, we instead measure \emph{top-K recall}---the
fraction of a reference method's top-$K$ molecules retrieved by the faster
method---and EF1\%, the hit enrichment in the top 1\% used by ROSHAMBO2. We compare against a patched version of ROSHAMBO2 fixing correctness bugs we identified (see Appendix~\ref{app:roshambo2fixes}).

Table~\ref{tab:main-results} reports the accuracy and wall-clock performance of ROSHAMBO3 as compared to ROSHAMBO2 on the same query sets. The field kernel retains ranking parity with ROSHAMBO3's exact analytic scoring (\FieldPanelRecallAtThousandPercent\% mean recall @ 1000, \FieldWorstRecallAtThousandPercent\% worst, with median pose displacement of
\FieldMedianPoseRmsdAngstrom~\AA). Notably, combo optimization often changed the winner and improved its objective, but not early
enrichment over rescore; meanwhile, combo took twice as long as shape-only screening whereas rescore paid only a 10\% penalty (Appendix Table~\ref{tab:public-performance}). ROSHAMBO3 therefore exposes all three modes.

ROSHAMBO3 was \VoltaShapeCorrectedSpeedup$\times$ faster than ROSHAMBO2 on shape-only search and \VoltaComboStockSpeedup$\times$ faster on combo search. For each mode we quote the harder of two baselines: ROSHAMBO2 as released, and ROSHAMBO2 with our color-skip performance fix (Appendix~\ref{app:roshambo2fixes}). Against the released version, the shape advantage rises to \VoltaShapeStockSpeedup$\times$.

\input{generated/tables/main-results.tex}

\paragraph{Large-scale search.} ROSHAMBO3 makes search over databases
of 10s-100s of millions of compounds feasible at near-interactive speeds.
On the RTX PRO 6000 Blackwell Max-Q, a PubChem3D screen aligns \DfiveConformers{}
conformers against \DfiveQueries{} queries in \DfiveWallSeconds~s end-to-end
(\DfiveKernelSeconds~s kernel), or \DfiveMillionEvaluationsPerSecond{} million
pair evaluations/s (Appendix Table~\ref{tab:scale-roofline}). For color combo on the internal library, increasing concurrent queries from 1 to 16 raises
throughput from \ComboQOneMillionEvaluationsPerSecond{} to
\ComboQSixteenMillionEvaluationsPerSecond{} million evaluations/s (Appendix Figure~\ref{fig:saturation}), and wall-clock time per query falls from
\ThreeCrQOneWallSecondsWhole~s to \ComboQSixteenWallSecondsPerQueryWhole~s. Concurrent queries therefore amortize the low-$q$ host limit until the device becomes limiting.

The field kernel reaches \RooflineBlackwellKernelToStructuralPercent\% of the  GB202 instruction issue bound and
\RooflineTitanKernelToStructuralPercent\% of the projected GV100 issue bound. The RTX PRO 6000 Blackwell consumes stored corpus bytes at
\PubchemStoredGigabytesPerSecond~GB/s when screening, demonstrating the importance of systems-level consideration of storage hierarchy. 


\paragraph{Portability.} The CPU and Metal backends share the CUDA backend's
corpus format, search, and ranking (Table~\ref{tab:portability},
Appendix~\ref{app:backends}). Every device's top-1,000 agrees with the CUDA
reference to at least \PortabilityLowestMeanRecallPercent\%
(\PortabilityWorstRecallPercent\% for the worst query); the reference GPU is
\PortabilityReferenceOverFastestCpuLow--\PortabilityReferenceOverFastestCpuHigh$\times$
faster than the fastest CPU, so CPUs extend ROSHAMBO3 to machines without NVIDIA
GPUs rather than replacing them.

\input{generated/tables/portability.tex}

\section{Conclusions}

\paragraph{Open-source shape similarity.} ROSHAMBO3 accelerates shape similarity by nearly two orders of magnitude at iso-accuracy and represents a new high-water mark for performance in open-source shape tooling. We anticipate that interactive-level performance on large databases will be useful for both human-driven and agentic chemical exploration. Nevertheless, opportunities remain, for example, identifying ways to use tensor units to add further speedup, and building a converged reference for the color objectives: our accuracy comparisons grade ROSHAMBO3 against another optimizer, and a far more exhaustive search for the color-rescore and combined shape--color optima would let any optimizer be graded against the best attainable score instead.

\clearpage
\paragraph{Agentic design of scientific software.} This work began with an expert question: how should a ``field'' shape method be
formulated for a modern GPU? Agents inspected unfamiliar code, implemented
isolated variants, constructed invariants, and interpreted profiles; they found
quaternion, optimizer, and color-skip defects in ROSHAMBO2, proposed the exponential map and Adagrad retuning, and iterated over hundreds of kernel parameterizations to explore design space.

\begin{figure}[t]
  \centering
  \includegraphics[width=\textwidth]{ticket-strip.pdf}
  \caption{The development record in brief (Appendix Figure~\ref{fig:ticket-timeline} draws it in full).
  One dot per ticket at its creation, colored by the tier of the implementing model, with a faint line to its closure;
  circles mark each epoch's main wave of ticket starts. Top: commits per day. Dashed lines mark the workshop's
  submission deadline, author notification, and camera-ready deadline.}
  \label{fig:ticket-strip}
\end{figure}

The collaboration changed as the project grew (Figure~\ref{fig:ticket-strip}; Appendix~\ref{app:agents}).
Parallel agents in separate worktrees built the GPU engine in
\AugustActiveDays{} active days. Once results began to reframe later questions,
the author dispatched one ticket at a time, each reviewed by a different model
tier; OpenAI Codex then joined the same repository and carried most of the CPU,
ARM, and Metal ports. Across \TicketCount{} tickets (\ClosedTicketCount{}
closed), \DecisionCount{} reopenable decisions, and \PitfallCount{} documented
failure modes, context lived in the repository, not in any session: each ticket
states its context, deciding measurement, and acceptance threshold up front,
closes only when that measurement is committed, and can be picked up by any
model or vendor's agent with no conversation. We think this is the most
reusable part of the workflow.

However, this work was not fully autonomous and expert direction remained essential: models proposed plausible but incorrect
bottlenecks and opposed experiments that later improved performance. As an example, the author rejected placing shape atoms in registers, even though it markedly improved the performance of shape optimization, intuiting that this would  become catastrophic for performance with color; later experiments confirmed this constraint. While the ticket-based mode of operation enabled effective local search by agents, long-horizon and long-context exploration still remained challenging, and regular feedback and context management were critical.

The relative ease with which intricate optimizations were performed by AI agents to unlock a 100$\times$
performance increase in a long-studied scientific kernel raises the question:
how many more such 10--100$\times$ improvements might be latent even in methods that have been worked on
for decades? No exceptional AI resources were needed to build ROSHAMBO3 --- one \ClaudePlan{} seat
(US\$\ClaudePlanUsdPerMonth/month) and one \ChatgptPlan{} seat (US\$\ChatgptPlanUsdPerMonth/month),
with modest additional credit usage ($<$US\$\AdditionalCreditUsdBelow). A token investment in tokens
by skilled tool developers may stand to produce a large compounding return for 
the broader scientific community.

\begin{ack}
Thanks to the authors of ROSHAMBO and ROSHAMBO2 --- Rasha Atwi, Ye Wang, Simone Sciabola, Adam Antoszewski, and Stephen Farr --- for resurrecting and improving
open-source GPU-accelerated shape-similarity software, and for graciously allowing me to put a new version number
on their lineage of tools.
\end{ack}

% References and appendices are outside the six-page camera-ready content limit.
\clearpage
\bibliographystyle{plainnat}
\bibliography{references}
\appendix
\newpage
\section*{Appendix}

\section{Tested ROSHAMBO2 configurations}
\label{app:roshambo2fixes}
At the outset of developing ROSHAMBO3, we identified a number of latent
bugs in ROSHAMBO2 affecting both accuracy and performance. Fixes for these
have been submitted upstream as pull requests \#10 (correctness) and \#11
(performance) to \url{https://github.com/molecularinformatics/roshambo2}, but
as of the time of writing, have not yet been merged. Briefly:
\begin{itemize}
	\item \textbf{Correctness}: Quaternion magnitude was incorrectly computed in normalization (typo of $*$ in place of $+$); in rare cases molecule preprocessing could mirror a compound inverting its stereochemistry; and a subtle bug mixed combo-optimization gradients with incorrect relative weights.
    \item \textbf{Performance}: with color scoring disabled, ROSHAMBO2 would still iterate over all color atom pairs but then post-multiply computed scores by zero. This computes a \emph{correct} shape-only score, but does much more work than is necessary.
\end{itemize}

In Table~\ref{tab:main-results}, the DUDE-Z enrichment comparisons use ROSHAMBO2 with the correctness fixes applied. The speed comparisons time two ROSHAMBO2 builds in one earlier session, as released and with only the color-skip performance fix, against ROSHAMBO3 in the release benchmark run on the same TITAN V host. The color-skip fix matters only when color scoring is off, so it speeds up the shape baseline but not combo. For each mode we lead with whichever baseline gives the smaller speedup: the color-skip build for shape and the released build for combo.

\section{Optimization mathematics and related work}
\label{app:methods}

This supplement is intended to make the numerical and systems choices
reconstructable, including experiments that did not land. The paper's build
regenerates every table, figure and constant from run records that ship in the
paper's repository, redacted where they concern private chemistry or
infrastructure, and its manifest records their hashes. The internal repository
contains confidential chemistry and infrastructure metadata and is not itself a
release artifact.

\subsection{Gaussian field and derivatives}

For a fixed query atom $i$, write
\begin{equation}
 \phi_i(\vect x)=C_i\exp[-\beta_i\|\vect x-\vect a_i\|^2].
\end{equation}
The exact spatial derivative sampled when the field is constructed is
\begin{equation}
 \nabla\phi_i(\vect x)=-2\beta_i(\vect x-\vect a_i)\phi_i(\vect x),
 \qquad \nabla\Phi_A=\sum_i\nabla\phi_i.
\end{equation}
If $\vect y_j=R\vect b_j+\vect t$, translation has derivative
$\sum_jw_j\nabla\Phi_A(\vect y_j)$. For an infinitesimal left rotation,
$\delta\vect y_j=\delta\vect\theta\times R\vect b_j$, giving the torque
\begin{equation}
 \frac{\partial V_{AB}}{\partial\vect\theta}
 =\sum_j w_j(R\vect b_j)\times\nabla\Phi_A(\vect y_j).
\end{equation}
Because self-overlaps are rigid-motion invariant, the derivative of the
Tanimoto score is a scalar multiple of the cross-overlap derivative:
\begin{equation}
 \nabla T=\frac{V_{AA}+V_{BB}}
 {(V_{AA}+V_{BB}-V_{AB})^2}\nabla V_{AB}.
\end{equation}

The implementation evaluates the analytic value and gradient at grid nodes in
fp32 and stores the four channels together in one fp16 texture element. CUDA
texture hardware performs the same trilinear interpolation coordinates for all
channels. Interpolating sampled analytic derivatives is not, in general,
identical to differentiating the piecewise-trilinear scalar interpolant; the
optimizer therefore follows a smooth approximation to the analytic gradient,
not a formally conservative derivative of the stored scalar field. We test the
consequence at the level that matters for screening---optimized poses and
retrieval---rather than concealing it behind value-only error.

The field is padded beyond the query support and uses zero border addressing.
This makes an excursion outside the allocated box approach the zero-overlap
limit safely. The production grid spacing was selected with a preregistered
retrieval gate over full spacing sweeps and per-query failures, not solely by
minimizing pointwise field error.

\subsection{Rigid-motion Adagrad}

Let $\vect g_{t,k}$ and $\vect g_{q,k}$ denote translation and tangent-space
rotation gradients. ROSHAMBO3 updates separate elementwise accumulators,
\begin{align}
 \vect G_{t,k}&=\vect G_{t,k-1}+\vect g_{t,k}^{\,2}, &
 \delta\vect t_k&=\eta_t\vect g_{t,k}/\sqrt{\vect G_{t,k}+\epsilon},\\
 \vect G_{q,k}&=\vect G_{q,k-1}+\vect g_{q,k}^{\,2}, &
 \delta\vect\theta_k&=\eta_q\vect g_{q,k}/\sqrt{\vect G_{q,k}+\epsilon}.
\end{align}
The quaternion update in Eq.~\ref{eq:expmap} preserves the group structure
without the singular behavior of additive quaternion coordinates. Separate
rates are physically meaningful: one controls translation in angstroms and the
other rotation in radians.

For an approximately stationary gradient magnitude, Adagrad step lengths decay
as $k^{-1/2}$. Thus $\eta_t\sum_{k=1}^{n}k^{-1/2}$ is the relevant translation
reach, rather than $\eta_t$ alone. The original tuning compared coupled and
uncoupled parameterizations across learning-rate and step-count sweeps. The
shape-derived budget was then independently retested for the combination
objective rather than inherited by assumption; the production budget remained
appropriate.

\subsection{Color modes and optimization semantics}

Color features are pseudo-atoms placed at pharmacophore locations rather than
heavy-atom volume centers. Let $A_u$ and $B_v$ contain the features of types
$u$ and $v$, where the six production types are hydrophobe, acceptor, aromatic,
donor, positive ionizable, and negative ionizable. The color cross-overlap is
of the form
\begin{equation}
 V^{\rm c}_{AB}(R,\vect t)=
 \sum_{u,v}\sum_{i\in A_u}\sum_{j\in B_v}
 K_{uv}\exp[-\tfrac12 r_{uv}
 \|\vect a_i-(R\vect b_j+\vect t)\|^2],
 \label{eq:color-overlap}
\end{equation}
where the type-pair interaction matrix supplies the amplitude $K_{uv}$ and
width $r_{uv}$; a zero entry suppresses an incompatible pair. The canonical
production matrix is diagonal, with common parameters for matching types, but
the implementation admits a general matrix. This is an implicit chemical
interaction model: feature typing and the matrix determine which spatial
coincidences are favorable.

We use \emph{plane} to mean one logical feature-type channel---the contiguous
feature list, or its corresponding scalar field---not a two-dimensional
geometric plane. A field implementation may therefore require one query field
per interacting source type. Shape and color self- and cross-overlaps are
computed separately to form a shape Tanimoto $T_{\rm s}$ and color Tanimoto
$T_{\rm c}$ before any weighted combination. The three public modes have
deliberately different semantics:

\begin{description}
 \item[Shape] optimizes and ranks only the shape Tanimoto.
 \item[Rescore] selects the best shape pose and evaluates color once at that
 pose; it preserves the shape search exactly.
 \item[Combo] mixes the separately normalized shape and color objectives inside
 every optimization step, so color may change the pose, winning start, and
 winning conformer.
\end{description}

Sparse color planes and the dense hydrophobe plane present different arithmetic
intensities. The development work consequently tested analytic, field, and
hybrid assignments by plane rather than assuming one representation was best
for every feature. We report the resulting implementation choices here and the
rescore-versus-combo quality comparison in Section~\ref{app:protocols}.

\subsection{Relationship to grid and Gaussian-sum methods}

The algebraic transformation in Eq.~\ref{eq:field} belongs to a broader family
of methods and should not be interpreted as the invention of either field-based
alignment or fast Gaussian summation. BRUTUS performs rigid molecular
superposition by optimizing a similarity between interpolated molecular
interaction fields~\citep{ronkko2006brutus}. Its objective is an interaction
field rather than Grant--Pickup first-order shape density, and it does not
describe evaluation of a fixed query density at each fit atom.

The closest mathematical precedent we found is the kernel-correlation method
of \citet{tsin2004correlation}. It represents point sets by Gaussian kernel
density estimates and notes that pair enumeration can be replaced by a sampled
fixed density, or ``affinity map.'' FilterReg later evaluates a fixed Gaussian
mixture at transformed points using fast filtering on a permutohedral
lattice~\citep{gao2019filterreg}. These methods address point-set registration
rather than molecular Tanimoto similarity, and use objectives and optimization
procedures different from those considered here. More generally, the Fast
Gauss Transform evaluates Gaussian sums in asymptotically linear
work~\citep{greengard1991fast}, but is neither a regular query grid nor a GPU
molecular-alignment implementation.

OpenEye's public Shape Toolkit exposes \texttt{Grid}, \texttt{Analytic}, and
\texttt{Exact} overlap methods and permits a reference grid to be compared with
a fit molecule~\citep{openeyeShapeTheory}. It is therefore plausible that ROCS
uses a related transformation. However, the public documentation does not
specify the field construction, interpolant, pose-gradient evaluation, GPU
mapping, or error characteristics needed to reproduce that implementation.

The contribution of \method{} is consequently the openly described
combination: a query field whose value at a fit atom is that atom's complete
first-order shape-overlap contribution; reuse of the field throughout a
library-scale search; simultaneous interpolation of value and pose gradient by
GPU texture hardware; integration with analytic color optimization; and an
end-to-end open implementation with measured retrieval agreement. To our
knowledge this complete combination has not previously been described in the
open molecular-shape literature. This is a claim about an operationally defined
implementation, not priority for the general idea of grids, fields, or fast
Gaussian evaluation.

\section{GPU and system engineering}

\paragraph{Execution model and occupancy.}
\label{sec:gpumodel}
An NVIDIA GPU comprises many streaming multiprocessors (SMs). CUDA assigns
thread blocks to SMs and divides each block into 32-thread \emph{warps}; the 32
threads, or lanes, execute a common instruction stream, while a warp scheduler
chooses a ready resident warp each issue cycle~\citep{nvidia2026cuda}. Keeping
several warps resident hides latency, but residency is finite: registers are
allocated per thread and shared memory per block. For example, our 256-thread
kernel at 64 registers/thread and four blocks/SM consumes the entire
65,536-register file. More live state either reduces resident blocks or spills
to ``local'' memory backed by device DRAM. Occupancy---resident warps divided by
the architectural maximum---is therefore a capacity measure, not a performance
objective: extra warps help only if the limiting execution or memory resource
can use them. This distinction is why we report schedule choices as measured
per-bucket outcomes rather than as a rule to maximize a single counter.

Modern CUDA GPUs feature a multi-stage cache hierarchy for global device memory, with L1 and L2 caches providing programmer-transparent caching of global memory accesses, and shared memory acting as a programmer-managed scratchpad. Textures use the same pool of L1 and L2 cache as other global memory accesses. Notably, L1 and shared memory divide the same pool of per-SM SRAM, creating a direct tradeoff for the field kernel between keeping data like coordinates resident in shared memory (avoiding a drop to L2 or global), versus leaving that L1 space available for field texture samples. Nsight profiles illustrate the resource shift the field produces:
% from the Nsight profile records; unchanged by this edit, only relocated.
analytic shape used \AnalyticShapeRegistersPerThread{} registers/thread and
kept only \AnalyticShapeActiveWarpsPercent\% of warps active, whereas the
\FieldShapeRegistersPerThread-register field path reached
\FieldShapeActiveWarpsPercent\%; the latter issued only
\FieldShapeFmaPercent\% of its instructions to the FMA pipe and ran L1TEX at
\FieldShapeLoneTexThroughputPercent\% and L2 at \FieldShapeLtwoThroughputPercent\%
of peak sustained throughput.

This design cannot be reduced to ``maximize occupancy.'' On the GB202,
shape-field work showed that doubling nominal occupancy could leave instruction
issue unchanged, while register-resident coordinates crossed from useful to
harmful as bucket width and spill traffic grew. A cache hint also reversed its
ranking for one bucket when the benchmark moved from the bucket mean to its
widest members. Adding color introduced state and reuse absent from shape and
made some shape-local register choices untenable. GV100 and GB202 then selected
different cells because their register, cache, instruction, and memory balances
differ. We report these reversals as results rather than smoothing them into one
architecture-independent rule; the internal record retains the complete sweep,
ptxas resource counts, profiler counters, and rejected cells.

\subsection{Shipped kernel configuration and explored design space}
\label{app:kernel-config}

The production field kernel has three user-visible modes: shape-only; shape
optimization followed by a terminal analytic color rescore; and combination
(``combo'') optimization, in which analytic color contributes to the objective
and gradient at every step. These modes share the field representation but do
not use one universal scheduling cell. Table~\ref{tab:shipped-kernel-config}
summarizes the resolved defaults for the standard four-start search. For a
one-start combo search, $S$ is clamped to one rather than leaving three of four
groups idle.

\input{generated/tables/shipped-kernel-config.tex}

\paragraph{Coordinate placement.}
Query shape has two simultaneous representations. The repeated optimization
samples a CUDA array of fp16 \texttt{half4} texels containing the field value
and its three spatial derivatives. The original query is also copied once per
block into shared memory as fp32 \texttt{float4}; this record supplies the
terminal exact shape score and any analytic color calculation. Query color
sites, their per-type segment bounds, and their Gaussian coefficients therefore
remain shared in both color modes. Color fields are implemented as an
experimental path but are disabled in the shipped configuration.

Library shape arrives in a compact three-array ($xyz$) int16 wire representation and
is staged once per library row into shared memory. All starts and all optimizer
steps then reuse that record. Terminal color rescore is too infrequent to merit
a second stage, so its library color coordinates are read from global memory.
Combo has much stronger reuse: the first 48 color sites of each library row are
staged as three int16 coordinate arrays, while the uncommon tail remains a
global-memory fallback. Color type is derived from the per-conformer count
vector rather than stored as a fourth plane.

\paragraph{Lane schedule and SRAM balance.}
$W_A$ is the number of cooperating lanes assigned to one candidate/start
alignment, while $S$ is the number of start groups resident in a warp row.
Increasing $W_A$ accelerates reductions within a pose but decreases independent
poses in flight. Increasing $S$ can recover pose-level parallelism, subject to
$S W_A\leq32$, shared-memory allocation, and register pressure. Atom-count
buckets make the best cell depend on the padded width as well as the real atom
count distribution inside that bucket.
Every production cell uses $W_A=4$, so four adjacent lanes divide one
candidate's atom loop. Shape and rescore use $S=1$: each four-lane group owns a
library row and loops over all starts, giving 64 rows per 256-thread block.
Four-start combo uses $S=4$: four groups share a row and each handles one start,
giving 16 rows per block. The shared shape record is consequently amortized
across the four groups. Component $c$ of the seven-component Adagrad update is
owned by lane $c\bmod4$; this reduces issued update instructions without
changing the atom reduction.

The preferred shared-memory carveout is explicitly set to 25\% on every
production launch, because the CUDA attribute is sticky across launches. On
both tested chips this request selected a 32-KiB physical partition, although
the percentage-to-partition mapping is architecture-specific. Shape and
rescore retain the four-block launch bound (a nominal 64-register/thread
budget). Combo uses a two-block bound and compiled at approximately 124--127
registers/thread; the extra state pays for greater issue activity rather than
maximum nominal occupancy.

\paragraph{Selection and negative results.}
The final point is a joint shape-and-color configuration, not the winner of a
shape-only microbenchmark. Register-resident library shape won an isolated
shape campaign, but the merged shape/color record enlarged its register tile;
every relevant color instantiation spilled and whole-screen throughput fell.
Ordinary global loads and the \texttt{\_\_ldcg}/\texttt{\_\_ldcs} cache hints also
failed to beat shared shape residency on both tested architectures. Staging 48
color sites helped GB202 but modestly hurt GV100; we nevertheless retain it as
the single portable default rather than dispatching by GPU model.

The measured search space included $W_A\in\{1,2,4,8,16\}$ and
$S\in\{1,2,4,8\}$; three update-lane arrangements; shared, global, cache-hinted,
and register coordinate tiers independently for shape, query color, and
library color; color stages from 16 to 64 sites; one to four launch-bound
blocks/SM; shared-memory carveout; six shape/color accumulator schedules;
several color pair-loop structures; field spacing and padding; texture formats
and boundary handling; fp32 versus reduced-precision records; and interleaved
versus planar library tiles. Several rankings reversed when color was added or
when the same cells moved between GV100 and GB202, which is why these choices
were gated by full-corpus experiments rather than inferred from occupancy or a
single profile counter.


\subsection{Operation accounting and a multi-resource roofline}
\label{app:roofline}

The ceilings in this section are for three specific chips and boards: the GB202
in our \RooflineBlackwellPowerLimitW-W RTX PRO 6000 Blackwell Max-Q, the AD104 in a \RooflineLFourPowerLimitW-W cloud L4, and the
GV100 in our \RooflineTitanPowerLimitW-W TITAN V. They are not ceilings for the Blackwell, Ada, and
Volta architecture families. In particular, GB200 and GB300 have substantially
different resource and power limits from this GB202 configuration. The three
also differ in what was measured: GB202 and AD104 were profiled with Nsight,
GV100 was not, and the L4 has no PubChem3D screen (its clock comes from its
LIT-PCBA run). We therefore use chip names for modeled
resources below and reserve product names for measured systems.

The production optimizer has fixed work: it performs 25 updates and one final
field evaluation for each of four starts, or
\RooflineFieldPassesPerPair{} field passes per query--candidate pair. There is
no convergence-based early exit. Let $A_0$ and $B_0$ be the real shape-atom
counts of a query and candidate, and let $A_t,B_t$ be their color-feature
counts by type. Defining $P=104$ and $C(A,B)=\sum_t A_tB_t$, the exact semantic
work per pair is
\begin{center}
\begin{tabular}{@{}lrrr@{}}
\toprule
mode & shape texture samples & exact shape terms & analytic color terms \\
\midrule
shape & $PB_0$ & $A_0B_0$ & 0 \\
rescore & $PB_0$ & $A_0B_0$ & $C(A,B)$ \\
combo, analytic color & $PB_0$ & $A_0B_0$ & $(P+1)C(A,B)$ \\
\bottomrule
\end{tabular}
\end{center}
The final terms are the exact analytic cascade at the winning pose. If a color
plane is field-served, its optimization term changes from $PA_tB_t$ Gaussian
evaluations to $PB_t$ texture samples; its exact final rescore is unchanged.
Query-field construction and self-overlaps are fixed per query and are
accounted separately from pair work.

\input{generated/tables/roofline-work.tex}

For PubChem3D, the manifest contains \DfiveLibraryRealAtoms{} real library atoms and
the four queries contain \DfiveQueryRealAtomsSum{} in aggregate. The measured
field kernel therefore sustains
\DfiveFieldBillionTextureSamplesPerSecond{} billion candidate-level
\texttt{float4} interpolations/s while its exact cascade evaluates
\DfiveFieldBillionExactGaussianTermsPerSecond{} billion Gaussian pairs/s. In
the model, transposition removes a factor of
\FieldOptimizerGaussianReductionFactor{} in optimizer Gaussian evaluations and
moves that work from CUDA-core and special-function instructions to the
texture path. The comparison is a movement on a resource-usage surface, not a
claim that a texture sample and an exponential are equivalent operations.

The matched explicit-sm120 GB202 profile confirms the predicted movement. The
analytic kernel executes \RooflineAnalyticToFieldInstructions{}$\times$ as many
warp instructions and \RooflineAnalyticToFieldXuInstructions{}$\times$ as many
XU instructions as the field kernel. It uses
\RooflineAnalyticRegistersPerThread{} rather than
\RooflineFieldRegistersPerThread{} registers/thread and sustains only
\RooflineAnalyticActiveWarpsPercent\% rather than
\RooflineFieldActiveWarpsPercent\% active warps. XU utilization falls from
\RooflineAnalyticXuUtilizationPercent\% to
\RooflineFieldXuUtilizationPercent\%, while L1TEX utilization rises from
\RooflineAnalyticLoneTexUtilizationPercent\% to
\RooflineFieldLoneTexUtilizationPercent\%, at
\RooflineFieldTextureSectorsPerRequest{} texture sectors/request. Thus the
field does not exchange one saturated resource for another: it reaches a more
balanced point with greater latency-hiding capacity. These are profiler-clock
diagnostics over identical pair rows, not wall-time measurements.

The three storage precisions do not imply three comparable arithmetic roofs.
Library coordinates cross H2D as int16 and are converted and scaled when loaded;
pose arithmetic then uses fp32. Query-field nodes are stored as fp16
\texttt{half4}, but texture hardware performs the filtered fetch and returns a
\texttt{float4}; there is no dense CUDA-core fp16 arithmetic loop to price
against nominal fp16 throughput. Pose transforms, gradients, reductions,
Adagrad state, and the exact cascade remain fp32.

Our issue count is therefore deliberately precision-agnostic: the executed
SASS count includes int16 loads and conversion, integer address/control work,
fp32 arithmetic, texture instructions, reductions, and optimizer updates. In
the matched field profile, scheduler issue reaches
\RooflineFieldIssueUtilizationPercent\%, while the general ALU and FMA pipes
reach \RooflineFieldAluUtilizationPercent\% and
\RooflineFieldFmaUtilizationPercent\%, respectively. Thus integer/conversion
work is included in the binding issue roof and is not close to an independent
ALU limit on GB202. The separate FP32 row is only a subresource ceiling, not a
claim that all kernel work is fp32. The GV100 issue result remains a projection,
not a measured per-pipe Volta profile. It uses the instruction demand profiled
on an sm\_89 build (an NVIDIA L4, weighted to PubChem3D the same way), because
the compiler emits the same field atom-loop shape for sm\_70 and sm\_89, whereas
sm\_120 unrolls it differently; per-pair counts remain specific to the build.

For semantic work we use a declared FLOP convention: add, subtract, and
multiply count as one; FMA counts as two; and \texttt{exp2}, square root,
reciprocal/division, sine, and cosine are reported separately. A source-level
model counts the dominant atom loops, while the SASS issue model includes the
compiler-dependent residual. Similarly, a half4 trilinear lookup has an
eight-node, 64-byte uncached stencil footprint, but this is not DRAM traffic:
stencil overlap, warp coalescing, border returns, and L1TEX/L2 reuse intervene.

Consequently a one-dimensional FLOP/DRAM-byte roofline is insufficient. For
evaluation rate $R$, the useful envelope is
\begin{equation}
 R \leq \min(R_{\rm disk},R_{\rm H2D},R_{\rm TEX},R_{\rm L1TEX},
 R_{\rm FP32},R_{\rm XU},R_{\rm issue},R_{\rm reduce},R_{\rm power}).
\end{equation}
For a device with $N$ SMs at observed clock $f$, $s$ warp schedulers, $c$
FP32 lanes, $x$ XU instructions/SM/cycle, $t$ texture units, and $l$ L1TEX
sectors/SM/cycle, we evaluate
\begin{align}
R_{\rm issue}&=Nsf/I, & R_{\rm FP32}&=2Ncf/F, &
R_{\rm XU}&=Nxf/X,\\
R_{\rm TEX}&=Ntf/T, & R_{\rm L1TEX}&=Nlf/S, &
R_{\rm DRAM}&=b_{\rm DRAM}/D,
\end{align}
where $I$ is the total executed warp-instruction count across all precisions and
pipes; $F,X,T,S,D$ are thread-level FP32 FLOPs, warp XU instructions, texture
requests, 32-byte texture sectors, and DRAM read bytes per pair evaluation.
These are independent upper bounds, not additive rates. We do not assign a
separate fp16-core roof because fp16 is a texture storage format here, nor an
integer-ALU roof because the captured ALU utilization is remote from issue;
either could be added if a future implementation moves substantial work there.

Each profile is a ladder of \RooflineLadderLaunches{} Nsight launches, one per
PubChem3D launch-width class, measured on GB202 and on AD104 (\RooflineLadderMinLibraryAtoms{} to
\RooflineLadderMaxLibraryAtoms{} real library atoms per candidate), each with the
same four queries as PubChem3D. Each launch's dynamic counts are weighted by its
class's share of PubChem3D pair evaluations, so no count is extrapolated beyond a
measured width; the mean is \RooflineDfiveMeanLibraryAtoms{} library atoms. On GB202 this gives
\RooflineFieldWarpInstructionsPerPair{} versus
\RooflineAnalyticWarpInstructionsPerPair{} warp instructions,
\RooflineFieldFpFlopsPerPair{} versus \RooflineAnalyticFpFlopsPerPair{} counted
FP32 FLOPs, and \RooflineFieldXuInstructionsPerPair{} versus
\RooflineAnalyticXuInstructionsPerPair{} XU instructions per PubChem3D-average pair.
The field path additionally generates \RooflineFieldTextureRequestsPerPair{}
hardware texture requests and \RooflineFieldTextureSectorsPerPair{} sectors per
pair. A warp request serves four candidate lanes, explaining why this request
count is one quarter of the candidate-level semantic lookup count. Occupancy
falls with launch width, so the field kernel's issue utilization runs from
\RooflineFieldIssueNarrowestPercent\% at the narrowest launch to
\RooflineFieldIssueWidestPercent\% at the widest; the percentages quoted below
are weighted by the time PubChem3D's pairs spend at each width. Per-pair
instruction counts are specific to the device and build that executed them.

Here ``fixed-clock'' means that the formulas use the median clocks observed in
production screens, \RooflineBlackwellObservedClockMHz{} MHz for GB202,
\RooflineLFourObservedClockMHz{} MHz for AD104 (its LIT-PCBA shape screen, under
the L4's power cap), and \RooflineTitanObservedClockMHz{} MHz for GV100, but still assume that every
listed subunit can sustain its peak per-cycle issue rate. It is not an
effective-throughput model for the device power limit. For example, the
GB202 warp-issue bound would be
\RooflineBlackwellMaxClockIssueMillionPairsPerSecond{} million pairs/s at the
\RooflineBlackwellSpecifiedMaxClockGHz-GHz specified maximum. Substituting the lower clock observed during the
\RooflineBlackwellPowerLimitW-W run reduces it to
\RooflineBlackwellFieldStructuralMillionPairsPerSecond{} million/s, but does
not account for instruction dependencies, unavailable warps, mixed-pipe
contention, or further activity-dependent power management. Our GB202 has \RooflineBlackwellSms{}
SMs, four warp schedulers and four texture units/SM, \RooflineBlackwellFpLanesPerSm{} FP32 lanes/SM, and
\RooflineBlackwellDramGigabytesPerSecond~GB/s specified DRAM bandwidth. AD104 (L4) has \RooflineLFourSms{} SMs, four
schedulers and texture units/SM, \RooflineLFourFpLanesPerSm{} FP32 lanes/SM, and \RooflineLFourDramGigabytesPerSecond~GB/s.
GV100 has \RooflineTitanSms{} SMs, four schedulers and texture
units/SM, \RooflineTitanFpLanesPerSm{} FP32 lanes/SM, and \RooflineTitanDramGigabytesPerSecond~GB/s. We use \RooflineXuWarpInstructionsPerSmCycle{} warp XU
instructions/SM/cycle on all three. The effective L1TEX rate is calibrated from each profiled
ladder (achieved texture sectors/SM/cycle over profiled L1TEX utilization):
\RooflineBlackwellLoneTexSectorsPerSmCycle{} sectors/SM/cycle on GB202
(\RooflineBlackwellLoneTexSectorsPerSmCycleMin{}--\RooflineBlackwellLoneTexSectorsPerSmCycleMax{}
across launches) and \RooflineLFourLoneTexSectorsPerSmCycle{} on AD104
(\RooflineLFourLoneTexSectorsPerSmCycleMin{}--\RooflineLFourLoneTexSectorsPerSmCycleMax{}). The GV100 projection
assumes \RooflineTitanLoneTexSectorsPerSmCycle{} and is explicitly a model rather than a TITAN V counter
measurement.

The field representation moves the structural bound away from XU issue. On GB202 the modeled analytic ceiling is
\RooflineBlackwellAnalyticStructuralMillionPairsPerSecond{} million PubChem3D-average
pairs/s, whereas the field ceiling is
\RooflineBlackwellFieldStructuralMillionPairsPerSecond{} million/s; on AD104
the corresponding values are
\RooflineLFourAnalyticStructuralMillionPairsPerSecond{} and
\RooflineLFourFieldLoneTexMillionPairsPerSecond{} million/s, and on GV100
\RooflineTitanAnalyticStructuralMillionPairsPerSecond{} and
\RooflineTitanFieldStructuralMillionPairsPerSecond{} million/s. On GB202 the field
kernel is bound by warp issue, with the L1TEX ceiling
(\RooflineBlackwellFieldLoneTexMillionPairsPerSecond{} million/s) close behind. On
AD104, whose counters are measured, L1TEX binds before issue
(\RooflineLFourFieldIssueMillionPairsPerSecond{} million/s); on GV100 the modeled
L1TEX ceiling also binds before issue
(\RooflineTitanFieldIssueMillionPairsPerSecond{} million/s), resting on the
assumed sector rate above rather than a Volta counter. Nominal texture-unit
issue and device DRAM are remote. This is the desired representation shift:
expensive analytic exponentials cease to dominate, and texture interpolation
uses a previously underused fixed-function path without saturating it first.

\input{generated/tables/roofline-device-caps.tex}

\input{generated/tables/roofline-observed.tex}

\input{generated/tables/scale-roofline.tex}

Power requires different treatment. Watts are not fungible operations, and
throughput is not globally linear in the power limit. The effective \RooflineBlackwellPowerLimitW-W
operating points are therefore the measurements in
Table~\ref{tab:roofline-observed}, not the bounds in
Table~\ref{tab:roofline-device-caps}. The RTX PRO 6000 Blackwell Max-Q ran at only
\RooflineBlackwellClockFractionPercent\% of its \RooflineBlackwellSpecifiedMaxClockGHz-GHz specified maximum,
drew \RooflineBlackwellMeanPowerW{} W on average with a
\RooflineBlackwellMaximumPowerW{}-W maximum sample, and reported software
power-cap throttling in \RooflineBlackwellPowerCapPercent\% of telemetry
samples. The observed kernel rate is
\RooflineBlackwellKernelMillionPairsPerSecond{} million pairs/s, or
\RooflineBlackwellKernelToStructuralPercent\% of the fixed-clock issue bound;
the \DfiveMillionEvaluationsPerSecond{}-million/s end-to-end rate is
\RooflineBlackwellWallToStructuralPercent\% of that bound because it also
contains pipeline time. A first-order scaling of the measured kernel point by
$300/\RooflineBlackwellMeanPowerW$ gives
\RooflineBlackwellLocalPowerExtrapolatedMillionPairsPerSecond{} million/s, of
which the kernel reaches \RooflineBlackwellKernelToPowerExtrapolatedPercent\%.
We report this only as a local sensitivity check: it is not an independent
roof or a prediction at a substantially different power limit.

The TITAN V, by contrast, ran at \RooflineTitanObservedClockMHz{} MHz
(\RooflineTitanClockFractionPercent\% of its \RooflineTitanSpecifiedMaxClockGHz-GHz maximum), averaged
\RooflineTitanMeanPowerW{} W, peaked at \RooflineTitanMaximumPowerW{} W under a
\RooflineTitanPowerLimitW-W limit, and recorded \RooflineTitanPowerCapPercent\% power-cap samples.
We therefore leave $R_{\rm power}$ inactive rather than extrapolating it. Its
\RooflineTitanKernelMillionPairsPerSecond{}-million/s matched PubChem3D kernel rate is
\RooflineTitanKernelToStructuralPercent\% of its fixed-clock issue bound; its
end-to-end rate reaches \RooflineTitanWallToStructuralPercent\%.
Finally, the isolated H2D and top-$k$ rates in
Table~\ref{tab:roofline-observed} are comfortably above each measured kernel
rate. They are not the binding stages for these runs, although disk production
and non-kernel wall time still lower end-to-end throughput.

Each ladder launch runs the production analytic and field kernels on the same
corpus slice for its launch-width class and records issue, pipe, instruction,
texture request/sector, cache, DRAM, occupancy, and register counters. Nsight
replays kernels and controls clocks, so profiled elapsed time is used only to
calibrate resource service rates and is never substituted for unprofiled
benchmark wall time. The release includes a single-launch profiling target
(\texttt{scripts/profile\_roofline.py}) and each ladder's counter summary and
pair weights, not the raw profiler reports; device parameters, calibration
inputs, tables, and constants regenerate from those summaries and the shipped
run records.

\subsection{Data representation, streaming, and storage}

The persistent corpus format quantizes coordinates on a fixed grid. Earlier
staging expanded each atom to a float4 record before transfer even though the
device did not require the fourth component. The production wire format sends
three int16 coordinates and expands them at consumption. An exhaustive int16
round-trip test precedes the optimization evidence. Where compiler scheduling
changed last-bit arithmetic, the project separated an encoding effect from a
binary-layout effect and required unchanged retrieval.

The PubChem3D bandwidth figures separate storage from device staging. Its
bit-packed coordinate sections occupy \PubchemStoredCoordinateGigabytes~GB and
the host decodes them to \PubchemDecodedCoordinateGigabytes~GB of padded int16
coordinates, a \PubchemCoordinateDecodeRatio$\times$ expansion; self-overlap
and run sections and the molecule table bring the stored corpus to
\PubchemStoredGigabytes~GB. The scan transfers
\PubchemHtwoDTransferredGigabytes~GB to the GPU,
\PubchemHtwoDToStoredByteRatio$\times$ the stored bytes: the decoded
coordinates plus per-conformer counts. Over the complete
\DfiveWallSeconds-s scan these are only \PubchemStoredGigabytesPerSecond~GB/s
of logical stored-corpus consumption and
\PubchemHtwoDWholeRunGigabytesPerSecond~GB/s of H2D bytes. The larger
\PubchemHtwoDActiveGigabytesPerSecond-GB/s figure divides H2D bytes by the
accumulated time inside copy calls. It is PCIe copy-engine bandwidth, not disk
bandwidth. Finally, this was a warm-cache run, so even the logical stored-byte
rate is not a measurement of physical NVMe traffic. Host-side decoding and
copy-time normalization explain the gap between the
\PubchemHtwoDActiveGigabytesPerSecond- and
\PubchemStoredGigabytesPerSecond-GB/s figures.

A format-3 corpus groups conformers into atom-count buckets, one
\texttt{.r3g} file per group, stores molecule/conformer identity beside
bit-packed coordinates and optional color features, and permits large
sequential reads. Host producers decode independent
chunks into pinned buffers. A single GPU-facing consumer preserves launch
ordering, overlaps transfers with earlier kernels, and retains only a bounded
molecule-level top-$k$. This distinction matters because conformer-level top-$k$
can return many poses of one molecule and because per-batch truncation is not
equivalent to a global reduction.

Both the 75M proprietary-library $q{=}1$ scan and the PubChem3D scan, on Blackwell,
consumed stored corpus bytes at a rate exceeding the capacity of 10-gigabit Ethernet
(\ThreeCrStoredGigabytesPerSecond~GB/s and \PubchemStoredGigabytesPerSecond~GB/s respectively,
compared to the \TenGigETheoreticalGigabytesPerSecond-GB/s theoretical bandwidth of 10-GbE before
protocol overhead). Thus, keeping a single Blackwell GPU fed from
remote storage would be a challenge without exceptional networking performance;
storage placement must be part of the benchmark configuration and production deployment design rather
than an incidental host detail.

\subsection{Scale and saturation}

\begin{figure}[H]
 \centering
 \includegraphics[width=5in]{blackwell-saturation.pdf}
 \caption{End-to-end combo throughput on the RTX PRO 6000 Blackwell Max-Q for
 the proprietary virtual hit-finding library.
 Points are medians of measured repeats. The last doubling yields only
 \ComboQEightToSixteenGainPercent\% additional throughput while runs reach the
 thermal or power cap.}
 \label{fig:saturation}
\end{figure}

\subsection{Systems and profile provenance}

Table~\ref{tab:systems} lists every benchmarked machine. Host, memory, GPU,
power-limit, and storage configurations bound end-to-end throughput, so they are
part of each benchmark rather than incidental details.

\input{generated/tables/systems.tex}
 Run records additionally preserve the
driver/toolkit, compiled architecture and fatbin targets, cache phase,
telemetry, and build identity; the shipped copies redact device serial numbers
and UUIDs. Warm-cache repeats characterize steady screening
service; cold-cache arms answer a different storage question and are never
silently pooled with them.

The internal record retains stage timings; register launch-bound, spill,
lane-assignment, cache-hint, and hydrophobe-field sweeps; wire-precision
ablations; and the causal sequence in which field acceleration exposed host
reduction before device top-$k$ removed dense D2H traffic. It also retains
profiler clocks and counter definitions so diagnostic profiles cannot be
mistaken for production wall times.

\section{CPU and Metal backends}
\label{app:backends}

\paragraph{Protocol.} Each device in Table~\ref{tab:portability} built the tool
from one clean source commit. It then screened the format-3 LIT-PCBA corpus
(the conformers of Section~\ref{sec:evaluation}, shape and color in one
build) with the eight PANEL-8 queries in shape,
rescore, and combo modes. Each mode had one warm-up and five timed repeats on
GPUs, three on CPUs, under each device's exclusive compute lock. By default,
the CPU backend runs one worker per logical CPU the process may use, SMT
siblings included. The Metal backend decodes the corpus on the CPU and holds
both locks. The cloud rows ran on dedicated, otherwise idle Google Cloud VMs.
On \ScalingIdenticalTopKHostCount{} of \ScalingHostCount{} hosts, every
placement returned an identical top-1,000.

Table~\ref{tab:portability}'s rate per two vCPUs (a core and its SMT sibling on
x86, two cores on Neoverse V2) is the full 16-vCPU VM's combo rate divided by
eight. Smaller placements on the same VMs ran somewhat faster per unit
(Table~\ref{tab:cpu-scaling}), so it is conservative for smaller VMs.
One driver script, published with the paper's data, runs every device. Its
per-host record names the install's source commit, the CPU kernel target
selected at run time, and the query file's hash.

\paragraph{Thread scaling.} Table~\ref{tab:cpu-scaling} screens one PANEL-8
query at each thread placement.
\begin{itemize}
\item \textbf{x86, one thread per core:} scaling stays close to linear up to
  eight cores, with combo efficiency of \ScalingXeonEightCoreEfficiencyPercent\% on
  the Xeon and \ScalingEpycEightCoreEfficiencyPercent\% on the EPYC.
\item \textbf{x86, the second SMT thread} adds combo throughput of
  \ScalingEpycSmtGainFourCorePercent\% on four cores and
  \ScalingEpycSmtGainPercent\% on all eight on the EPYC 9B45, and
  \ScalingXeonSmtGainFourCorePercent\% and \ScalingXeonSmtGainPercent\% on
  the Xeon.
\item \textbf{Neoverse V2} has no SMT. Its combo efficiency is
  \ScalingNeoverseEightCoreEfficiencyPercent\% on 8 cores and
  \ScalingNeoverseAllEfficiencyPercent\% on all 16.
\item \textbf{M4 Pro:} macOS places the threads, so the split between
	performance (P) and efficiency (E) cores is inferred. Shape screening runs at about
  \ScalingMFourProCpuShapePerPcoreThousand{}K pairs/s per P-core and
  \ScalingMFourProCpuShapePerEcoreThousand{}K per E-core. The four
  E-cores add \ScalingMFourProCpuEcoreGainPercent\% to combo
  throughput (10+4 against 10+0).
\end{itemize}

\input{generated/tables/cpu-scaling.tex}

\section{Experimental protocols}
\label{app:protocols}

\subsection{Corpus and query construction}

The LIT-PCBA source is the union of all 15 published assay sets. Canonical-SMILES
deduplication produced \LitpcbaSourceCompounds{} inputs. With nvMolKit \LitpcbaNvmolkitVersion{} and RDKit \LitpcbaRdkitVersion{},
we requested \LitpcbaConformersRequested{} ETKDGv3 conformers (seed \LitpcbaConformerSeed{}) and minimized each with MMFF94s;
\LitpcbaCompoundsWithoutConformers{} molecules produced no conformer, leaving \LitpcbaCorpusMolecules{} molecules
and \LitpcbaCorpusConformers{} conformers. Of the successful molecules, \LitpcbaFullConformerMolecules{}
reached all \LitpcbaConformersRequested{} requested conformers. The uncompressed conformer SDF was \LitpcbaSourceSdfGigabytes~GB;
\texttt{roshambo3 prepare} converted it to the
\LitpcbaDiskGigabytes-GB color-capable format-3 corpus.

PANEL-8 was selected before timing: sort the 15 target/assay groups, take the
first eight, and within each group take the first deposited template record.
The resulting queries are ADRB2/3p0g (27 heavy atoms), ALDH1/4wp7 (32),
ESR1-agonist/1l2i (24), ESR1-antagonist/1xp1 (34), FEN1/5fv7 (23), GBA/2v3d
(15), IDH1/4i3k (17), and KAT2A/5h86 (53). No embedding or minimization was
applied to these queries: timing uses their deposited coordinates. The field
accuracy panel is separate and broader (32 queries), preventing a hardware
timing panel from becoming the approximation gate.

PubChem3D is not a random or favorable subset; it is a complete copy of the 10-conformer-per-molecule PubChem3D
database: every available CID and up to ten deposited
conformers, giving \DfiveMolecules{} molecules and \DfiveConformers{}
conformers. The proprietary
virtual hit-finding library and its query are private. For its saturation sweep,
$q=1, 2, 4, 8,$ and 16 mean that
many conformers of the same \ThreeCrQueryHeavyAtoms-heavy-atom query, generated
by ETKDGv3 with seed \texttt{0xC0FFEE} and MMFF optimization; q is therefore a
controlled reuse factor, not a change of chemical workload.

\input{generated/tables/public-performance.tex}

The public LIT-PCBA arm is intentionally smaller than the PubChem3D headline but is
straightforward to rebuild from public sources
(\texttt{scripts/build\_litpcba.py}). PubChem3D supplies the public-data scale
result. The proprietary-library arm is retained only
for internally reproducible color saturation and mode-cost measurements.

\section{Agent-assisted development}
\label{app:agents}

This appendix describes how ROSHAMBO3 was built: how the collaboration between
the author and AI coding agents changed over two months, the ticket system that
carried context between sessions, models, and vendors, and the practices and
failure modes worth carrying to other scientific-software projects. Counts are
generated from a sanitized export of the development record (ticket ids,
epochs, model tiers, dates, and dependencies), which ships with the paper's
repository: \TicketCount{} tickets (\ClosedTicketCount{} closed) joined by
\TicketDependencyCount{} dependency edges, \DecisionCount{} standing decisions,
\PitfallCount{} numbered pitfalls, and \DevelopmentCommitCount{} commits over
\DevelopmentActiveDays{} active days. Figure~\ref{fig:ticket-timeline} draws
the whole record.

\subsection{How the collaboration evolved}

\paragraph{Before the repository.} The design question---how a ``field'' shape
method should be formulated for a modern GPU, and whether newer hardware units
could help---was explored first in chat conversations with Google's Gemini 3.0
Pro Preview. Nothing from this phase was checked in; its product was the
author's conviction that the field formulation was worth building.

\begin{figure}[t]
 \centering
 \includegraphics[width=\textwidth]{ticket-timeline.pdf}
 \caption{The development record. Each row is one ticket in id order, drawn
 from creation to the commit that closed it and colored by the tier of the
 model that implemented it; faint bars are tickets still open at export. Ids
 are assigned in creation order, so the record reads as a staircase. Shaded
 bands group tickets by epoch. Some bars end together where tickets were closed
 in a batch update of the ticket records (early September, and 10 October), not
 where their work ended. Dashed lines mark the workshop's submission deadline,
 author notification, and camera-ready deadline.
 Top: commits per day across the development repository, the ROSHAMBO2 fork where the ticket system
 began, and the separate workspace where the Metal port was done.}
 \label{fig:ticket-timeline}
\end{figure}

\paragraph{Plan and parallel swarm (15--29 August).} Work began in a fork of
ROSHAMBO2. A Claude Code agent wrote a planning document in which every
quantitative statement was tagged as either measured or modeled, and the same
evening decomposed it into 22 tickets with a dependency graph, a wave table of
what could run concurrently, and a model tier per ticket. Agents then worked in
parallel, each in its own git worktree, and the fork accumulated
\OriginCommitCount{} commits in five days. On 20 August the ticket system,
decisions, and results moved into a clean ROSHAMBO3 repository, and for the
workshop deadline six to eight lanes ran concurrently across two GPU machines.
The GPU engine of the submitted paper was built in \AugustActiveDays{} active
days.

\paragraph{Serial dispatch and review (from 2 September).} After submission the
author stopped dispatching waves and moved to one ticket at a time, each
started only on explicit approval, with an instruction to report the deciding
measurement and stop. The reason, recorded at the time, was that tickets in a
ladder were coupled: one result often invalidated the framing of the next
ticket, and running ahead spent effort on questions already answered. The same
day, tickets gained a \texttt{reviewer} field, and implementation and review
were assigned to different model tiers. The CPU port began in this mode.

\paragraph{Two agent systems (from mid-September).} OpenAI's Codex agents, first
used in August for a corpus build (T82), worked the repository alongside
Anthropic's Claude Code from mid-September. The two systems shared no
conversation; everything one needed from the other was in the files---tickets,
decisions, conventions, and results. By October Codex carried most of the
implementation: the CPU color and combo paths, the native x86 and ARM ports
(10--13 October), and, once the ARM port existed, the macOS and Metal port,
which took a single day.

\paragraph{The author's role.} Over the project the author's prompts moved from
specifying implementations toward three activities. \emph{Gating}: approving
the next ticket, closing one, or dispatching a review. \emph{Short expert
corrections} that re-aimed a line of work, often in a sentence: a field build
under 100~ms is irrelevant once amortized over a library; optimize headline
pairs/s, not relative savings; do not change the on-disk format while the GPU
cannot be benchmarked. And \emph{budget and drift control}: pausing reviewers to
save usage, or asking why convergence analysis had begun an hour into a
day-long job. Agents expanded implementation bandwidth; the author supplied
priorities, long-horizon constraints, and skepticism.

\subsection{Tickets as durable context}

Agent sessions are short, context windows are finite, and over this project the
models themselves changed several times. The repository is the only thing that
persists, so the ticket system was designed to put a unit of work into it in a
form that any agent---``possibly on a smaller model,'' in the words of its
README---could start without a conversation or the planning document.

%\clearpage
\paragraph{What a ticket contains.} Each ticket is one Markdown file with
structured frontmatter (id, lane, status, model tier, reviewer tier,
dependencies, branch, worktree, and artifact path) and six sections:
\emph{Context}, self-contained in a few sentences; \emph{Scope}, carrying the
plan's technical detail verbatim with its measured-or-modeled tags;
\emph{Deciding measurement}, the number that closes the ticket and where it must
be measured; \emph{Acceptance}, the thresholds, stated before work begins;
\emph{Traps}, the numbered pitfalls that apply, restated in the ticket's own
terms; and \emph{Working rules}, cited rather than copied, because rules copied
into 22 tickets had drifted apart.

\paragraph{Lifecycle.} A ticket is \emph{open}, \emph{claimed} (a worktree
exists and the claim is committed to the main branch), \emph{in review} (the
deciding measurement is committed under the ticket's results directory), or
\emph{closed}. A ticket closes only when its deciding measurement exists in the
repository, never on an agent's report of it; a negative answer closes a
ticket exactly as a positive one does. The in-review state exists because
agents do report success that their evidence does not support: the agent that
closed one ticket later recorded, in the ticket itself, that it had closed the
ticket on its own report and that review had been right to reopen it.

\paragraph{One source for each fact.} Status lives only in the ticket's
frontmatter and dependencies only in the dependent ticket; the index, the
reverse edges, and the graph checks (targets exist, no self-edges, no cycles)
are generated. Each rule followed a measured failure. When the ROSHAMBO2 fork
kept tickets in \texttt{open/} and \texttt{closed/} directories \emph{and} in a
status field, the two disagreed \StatusDriftCount{} times. When dependencies
were stored on both ends, \OneSidedBlocksEdges{} forward and
\OneSidedDependsEdges{} reverse edges lost their partners, and nothing noticed
because nothing read them.

\paragraph{Decisions and pitfalls.} Conclusions that outlive a ticket became
decisions, each with the measurement behind it and an explicit condition for
reopening it. They prevent settled questions from being relitigated by each
new session, and they are reversible by design. One recorded that the
exponential-map rotation update was not worth pursuing; the author reopened it,
and the entry was marked ``withdrawn, same day'' when the ticket that had
introduced the update showed a \ExpMapRecallGain{} recall gain once a covering
start strategy was removed. Another withdrew an apparent Volta speed inversion
as an artifact of a probe harness, and recorded the lesson as its own decision,
``a probe is not a kernel''.
Recurring failures became numbered pitfalls, appended in the order they were
discovered and cited by number in tickets (Table~\ref{tab:pitfall-classes}).
The repository's agent instructions direct every session to read the decisions
and pitfalls before starting work, and per-project memory notes carried the
author's standing preferences, such as serial dispatch, between sessions.

\input{generated/tables/pitfall-classes.tex}

\subsection{Several models, two vendors}

Each ticket names the tier of model expected to implement it and, from
September, a reviewer of a different tier; \ReviewedTicketCount{} tickets carry
a reviewer. Review was mostly across tiers of one vendor: Opus-tier work
reviewed by Fable, or a Luna-tier Codex implementation reviewed by Sol. Review
across vendors was rarer. Early in the project Gemini 3.1 Pro Preview reviewed
Opus's work, and agents of one vendor incidentally read, tested, and corrected
code written by the other's in the course of their own tickets. Reviews earned their
cost. One found a live bug in which a refactor of four field-gradient variants
had dropped a compensating rescale, leaving their torque at \ReviewTorqueScale{}
of its correct value. Another found that a tolerance derived for a
cross-implementation comparison was \ReviewToleranceLooseLow--\ReviewToleranceLooseHigh$\times$
looser than the measured cross-implementation distribution.

The two systems name their tiers differently (Fable, Opus, Sonnet, and Haiku
correspond to Astra, Sol, Terra, and Luna), and both vocabularies remain in the
ticket files. An attempt to normalize them to one scheme was reverted within
minutes, because each system reads the files and the other's names would have
broken its dispatch; the generated index renders both, so the mapping is read
rather than remembered. The models changed during the project:
Claude Sonnet and Opus from 4.8 through 5.5 and Fable 5 in Claude Code, and in
Codex GPT~5.6-Terra and Sol, then Sol, Astra, and Luna 6, and Sol and Luna 6.1.
Because the ticket carries the context, a change of model or vendor changed who
did the next ticket, not what the next ticket said.

\subsection{Practices worth reusing}

Each practice below exists because its absence cost the project something
measurable.
\begin{itemize}
 \item \textbf{A check that does not run cannot fire.} One ticket specified a
 stop rule in prose; no script evaluated it, and a stage was launched at
 \StopRulePlannedOverrunFactor$\times$ its specified size and had run
 \StopRuleExecutedOverrunFactor$\times$, \StopRuleOverrunHours{} hours of
 two-GPU time, before a person looked. Earlier instances were a
 refusal path that could not be reached, a test fixture that skipped instead of
 failing, and never-called grading rules. We make stop rules and
 acceptance thresholds code that writes its value into the run record.
 \item \textbf{Store each fact once and generate every view of it.} Helps avoid
 status / dependency drift above.
 \item \textbf{Correctness before timing, and the right kind of correctness.}
 Changes that only move data must be bitwise identical; changes to arithmetic
 must preserve retrieval (top-$K$ recall) against a reference. Neither gate is
 waived for a speedup.
 \item \textbf{Make a measurement quotable or call it advisory.} Any number
 cited in a decision is taken under an exclusive filesystem lock on the GPU, with
 clock, power, cache state, build identity, and repeat phase recorded. A
 profiler silently pins clocks, so numbers taken at different clocks are never
 combined into one ratio.
 \item \textbf{Assert what was loaded.} From a worktree, a shared Python
 environment silently imported the main checkout's code; a stale compiled
 extension failed exactly like a code bug. Runs record and check the source and
 build they executed.
 \item \textbf{Commit before waiting, and wait on a condition.} One ticket
 accumulated \WaiterShellCount{} concurrent polling shells with blind sleeps of
 up to \BlindSleepMaxSeconds~s while its result already sat in a log.
 \item \textbf{Siblings are not intruders.} Agents working in parallel see each
 other's worktrees, locks, and processes; one nearly deleted three live sibling
 worktrees as debris. The rule is to report and stop, never to clean up another
 agent's state.
 \item \textbf{Name code for what it does.} Identifiers named after tickets read
 as meaningful to the agent that wrote them and as noise to the next; the name
 says what code does, and a comment says where it came from.
 \item \textbf{Dispatch serially once results depend on one another, and bound
 scope in writing.} Parallel waves suit independent questions. When each answer
 reframes the next question, one ticket at a time with a written statement of
 what is \emph{not} authorized wastes less work.
\end{itemize}

\subsection{Where agents went wrong, and where the expert mattered}

The workflow was strongest at reconnaissance of unfamiliar code, proposing
independent implementations, constructing adversarial tests, and keeping
hypotheses alive, but had characteristic failures:

\paragraph{Plausible stories adopted before measurement.} Three tickets
attributed a kernel's cache pressure to library-molecule traffic; a profile
showed \QueryShareOfTexRequestsPercent\% of the relevant load requests came from
the query. A CPU scaling anomaly was explained as cold cache, which was
impossible for a \ImplausibleWorkingSetKB-KB working set against a
\ImplausibleCacheMB-MB cache; review found threads stacked on the same cores. A
broken control once produced an apparent \BrokenControlSpeedup$\times$ win that
was \CorrectedControlSpeedup$\times$ once corrected.

\paragraph{Opposing what later helped, or advocating what later hurt.} An
agent's decision ruled out the exponential map, which the author reopened and
which shipped. Conversely, agents favored keeping library shape coordinates in
registers, which cut whole-screen kernel time \RegisterShapeSpeedup$\times$ on a
shape-only screen; the author blocked it on the expectation that color would make it
untenable, and with color it measured \RegisterColorLow--\RegisterColorHigh$\times$.

\paragraph{Drift over long horizons.} A published head-to-head ratio
(\StaleHeadlineSpeedup$\times$) was still being cited after
\StaleHeadlineLandedTickets{} tickets had changed the engine underneath it, and
which ROSHAMBO2 build was the baseline was re-derived, inconsistently, by each
ticket that needed it until a decision fixed it. Long-horizon coherence came
from the author and from writing conclusions down, not single-session
memory.

\subsection{Limits of the record}

The tier in a ticket names a capability level, not a vendor, and git commit
trailers identify the orchestrating session rather than the model that wrote
each commit, so the record cannot attribute every ticket to a specific model.
Tickets for the Metal port were written after the fact from a separate
workspace (\RetroactiveTicketCount{} tickets). Some closures were batch updates
of the ticket records and do not date the end of the work. The paper's repository
ships the export, the conventions, and an excerpt of tickets and decisions; the
full archive, its results and its history hold confidential detail and are not
released. The record is the project's own: it shows what this collaboration
did, not what another team or a fully autonomous system would have done.

\end{document}
